\section{Rank coverage, calibration variability, and the specified shift}
\label{app:rank}
Condition on proper-training data and all independently chosen training randomness. Assume the resulting scalar calibration/new scores $S_1,\ldots,S_n,S_{n+1}$ are exchangeable, and use $k$ and $q$ from Equation~\eqref{eq:rank}. Break ties by independent random ordering for this argument only. The new score has uniform rank among $n+1$ scores. If its rank is at most $k\leq n$, then $S_{n+1}\leq S_{(k)}$ among the calibration scores. Therefore
\begin{equation}
 \Pr\{Y_{n+1}\in C_\alpha(X_{n+1})\mid\text{training}\}
 \geq \frac{k}{n+1}\geq1-\alpha.
 \label{eq:rank-proof}
\end{equation}
For distinct scores the first inequality is equality. Including ties with $\leq$ can only increase coverage. If $k=n+1$, the infinite threshold gives coverage one. Intersecting with the simplex preserves this event because the target lies there. A fixed positive training-derived scale and deterministic transformation in Equation~\eqref{eq:smoothing} preserve score exchangeability. The proof averages over calibration and a new observation; it does not condition on a realized calibration sample or establish pointwise conditional coverage.

For each coordinate $j$, calibration at $\alpha/K$ gives marginal coordinate miscoverage at most $\alpha/K$. The union bound then bounds any-coordinate miscoverage by $\alpha$, with no coordinate-independence assumption. Calibrating each at $\alpha$ instead gives only the bound $K\alpha$; it supplies no joint $1-\alpha$ guarantee.

For fixed observable group $g$, condition also on its realized calibration count $n_g$. Exchangeability of within-group scores yields the same bound for a new group-$g$ observation using its corrected rank, including infinity for insufficient calibration ranks. In the specified stress arm, the law of $(X,Y)\mid G=g$ is unchanged; only the test weights $w_g=\Pr_{\rm test}(G=g)$ change. Thus the group bounds imply
\begin{equation}
 \Pr_{\rm test}(Y\in C(X)\mid\text{training},\{n_g\})
 =\sum_g w_g\Pr_{\rm test}(Y\in C(X)\mid G=g,\text{training},\{n_g\})
 \geq 1-\alpha.
 \label{eq:shift-proof}
\end{equation}
The probability still averages over within-group calibration values and the future observation. This protection for fixed-group calibration does not extend to an arbitrary change inside groups. The stress arm does not evaluate density-ratio-weighted conformal prediction or general conditional-response shift.

\paragraph{Continuous-score diagnostic.}
Under the stronger assumption of conditionally IID scores with continuous CDF $F$, the transformed scores are IID uniform. For $k\leq n$, true coverage conditional on the realized threshold is $P=F(S_{(k)})$, hence
\begin{equation}
 P\sim\mathrm{Beta}(k,n+1-k),\quad
 \mathbb E[P]=\frac{k}{n+1},\quad
 \operatorname{Var}(P)=\frac{k(n+1-k)}{(n+1)^2(n+2)}.
 \label{eq:beta}
\end{equation}
This is the uniform order-statistic distribution, obtained by counting whether at least $k$ of $n$ uniform scores lie below a given value. For $M$ conditionally independent tests, the covered count conditional on $P$ is binomial. Total variance gives
\begin{equation}
 \operatorname{Var}(\widehat c)=\operatorname{Var}(P)+\frac{\mathbb E[P(1-P)]}{M}.
 \label{eq:beta-binomial}
\end{equation}
The saved uniform-score simulation and Figure~\ref{fig:calibration-diagnostic} check these references. IID continuity is required for this Beta law; exchangeability alone suffices for Equation~\eqref{eq:rank-proof} but not this distributional identity.

\section{Exact generator and fitting specification}
\label{app:generators}
Let $j=0,\ldots,K-1$, $\phi(X)=(1,u,v,\sin(\pi u),uv,G)$, and let the component coefficient vector be
\begin{equation}
 a_j=\left(.4\cos j,\ .7\cos\frac{2\pi j}{K},\ .5\sin\frac{2\pi j}{K},\ .3\cos\frac{4\pi j}{K},\ .2\sin\frac{4\pi j}{K},\ .2\cos\frac{2\pi j}{K}\right)^T.
 \label{eq:generator-coef}
\end{equation}
Write $\eta_j=\phi a_j$ and $\mu=\operatorname{softmax}(\eta)$. Interior, heterogeneous and zero mechanisms use this center. Boundary and Dirichlet mechanisms first replace $\eta_0$ by $\eta_0+3$. For logistic-normal mechanisms,
\begin{equation}
 Y=\operatorname{softmax}\{B[z(\mu)+\sigma_G Z]\},\qquad Z\sim N(0,I_{K-1}),
 \label{eq:logistic-normal}
\end{equation}
with the scales in Section~\ref{sec:methods}. All observations use fresh independent generator draws. For the Dirichlet mechanism set $\mu_D=.8\mu+.2\mathbf1/K$, sample independent $W_j\sim\operatorname{Gamma}(15\mu_{Dj},1)$, and take $Y_j=W_j/\sum_lW_l$. For the zero mechanism, independently with probability $.25$ select $J$ uniformly from $\{0,\ldots,K-1\}$, set the logistic-normal vector's $J$th entry to zero and divide by its remaining sum. The zero/positive slices condition on the realized response; they are descriptive audits rather than test-observable calibration groups.

Let $\Phi$ contain the first 1,500 training feature rows and $Z_{\rm train}$ their ILR response rows, with fixed zero smoothing where stated. The location coefficient matrix and inverse transform are
\begin{equation}
 \begin{aligned}
 \widehat A&=(\Phi^T\Phi+\Lambda)^{-1}\Phi^TZ_{\rm train},\quad
 \Lambda=\operatorname{diag}(0,.001,.001,.001,.001,.001),\\
 m(x)&=\operatorname{softmax}\{B\widehat A^T\phi(x)^T\}.
 \end{aligned}
 \label{eq:ridge}
\end{equation}
On the remaining 500 training observations, take the RMS of ILR residual norms within each group, using overall RMS for an absent group and flooring at $10^{-8}$. These observations do not fit the location. The predictor is intentionally shared across geometries and is misspecified for some mechanisms; conformal validity does not require correct specification. Generator/training implementations are \texttt{analysis/src/compbench/data.py}; the orthonormal basis and transforms are in \texttt{geometry.py}. The frozen manifests retain all trial seeds, split sizes and sensitivity choices.

\section{Common-simplex integration and empirical numerical allowance}
\label{app:integration}
Let $d=K-1$, $V_d=\pi^{d/2}/\Gamma(d/2+1)$, and $\operatorname{vol}_d(\simplex)=\sqrt K/d!$. These formulas use Hausdorff measure in the affine hyperplane. The primary evaluation integrates only $K=3$ cases.

\paragraph{Euclidean and coordinate regions.}
A tangent-space Euclidean ball of radius $r$ has volume $V_dr^d$. Its normalized in-simplex area is
\begin{equation}
 a_{\rm L2}=\frac{d!}{\sqrt K}V_dr^d\,\mathbb E_U[\mathbb1\{m+rBU\geq0\}],\qquad U\sim\mathrm{Uniform}(\text{unit }d\text{-ball}).
 \label{eq:l2-integral}
\end{equation}
For coordinate boxes, integrate the first $d$ barycentric coordinates in their component bounds; the last is one minus their sum. The coordinate-chart Hausdorff factor $\sqrt K$ divided by simplex volume leaves factor $d!$. For TV, $\tfrac12\|y-m\|_1\leq r$ implies $|y_j-m_j|\leq r$ because positive and negative differences each sum in magnitude to TV. This supplies a bounding box in the first $d$ coordinates; retain exact simplex/TV membership inside it. In $K=3$, an untruncated interior TV ball has normalized area $6r^2$, providing an analytic implementation check.

\paragraph{ILR and smoothing.}
For $w=\operatorname{softmax}(Bz)$, the tangent-space Jacobian is $K\prod_j w_j$. To see this, the derivative in the tangent space is represented by $B^T(\operatorname{diag}w-ww^T)B$; its determinant is the product of the nonzero eigenvalues of $\operatorname{diag}w-ww^T$. Each principal cofactor of that matrix is $\prod_j w_j$, by the determinant lemma, and the sum of the $K$ cofactors is that eigenvalue product. Hence, for $Z$ uniform on the radius-$r$ ball around $z(m)$,
\begin{equation}
 a_{\rm ILR}=\sqrt K\,d!V_dr^d\,\mathbb E_Z\!\left[\prod_j\operatorname{softmax}(BZ)_j\right].
 \label{eq:ilr-integral}
\end{equation}
With smoothing, center the ball at $z(T_\epsilon(m))$, write $w=\operatorname{softmax}(BZ)$, multiply the integrand by $(1+K\epsilon)^d$, and retain only $w_j\geq\epsilon/(1+K\epsilon)$ for every $j$. This follows from $y=(1+K\epsilon)w-\epsilon\mathbf1$ and computes the original target region's area, including its clipping at the simplex boundary. Boundary faces themselves have zero full-dimensional area.

\paragraph{Hellinger.}
Set $s=\sqrt y$ and $a=\sqrt m$, so both lie on the unit sphere in the positive orthant. The score constraint is $s^Ta\geq1-r^2$. The map $s\mapsto y=s^2$ has surface-to-simplex Jacobian $2^d\sqrt K\prod_j s_j$. Indeed, its squared tangent Gram determinant is $4^dK\prod_j s_j^2$: apply the cofactor identity to the diagonal squared derivative and use $\sum_js_j^2=1$. Division by simplex volume gives normalized Jacobian $2^dd!\prod_js_j$. Integrate this over the spherical cap, rejecting negative $s_j$. For $0<r<1$, cap sampling uses $\sin^2\theta\sim\mathrm{Beta}(d/2,1/2)$ truncated at $1-(1-r^2)^2$. If $r\geq1$ the entire simplex is covered. Infinite thresholds give area exactly one for every implemented family; finite radii can also cover the whole simplex in some geometries.

\paragraph{Paired refinement and allowance.}
The audit uses confirmatory repeats $0,25,50,75$ in every $K=3$ cell and the first four selected locations per group. At a matched location for methods $A,B$, let $D_{\rm main}$ be the main-resolution area difference and $D_s$, $s=1,\ldots,4$, the higher-resolution differences, pairing scrambles within each method pair. Define
\begin{equation}
 e=\max\left\{.001,\left|\overline D-D_{\rm main}\right|+t_{.975,3}\frac{\operatorname{SD}(D_1,\ldots,D_4)}{2}\right\}.
 \label{eq:allowance}
\end{equation}
Take the maximum over audited locations/repeats for each cell/group/pair, then the maximum over the two groups for its overall comparison. Enlarge each paired outer-trial interval by this amount on both sides. This empirical sensitivity policy can expose unresolved numerical orderings; its finite audit does not certify all unaudited sets or provide a simultaneous confidence bound. The main outer-trial interval remains a separate statistical uncertainty summary. The exact allowance implementation is \texttt{analysis/scripts/numerical\_allowance.py}.
